\vfill%
\pagebreak
\section{Iteration loop}%
\label{sec:flow:for_loop}
\label{sec:flow:for_loops}

\token{for} evaluates its body once for each element of an iterable value, with
one or more \textit{iterators} -- variables declared by the loop -- holding the
current element. Ranges, arrays, slices, tuples and maps are iterable
(Chapter~\ref{chap:compound}), and so is a class that implements the iteration
operator (Section~\ref{sec:class_override_for_loop}).

\subsection{Grammar}

\begin{lstlisting}[style=grammarVerb]
for_expression := 'cte'? 'for' iterator (',' iterator)*
                  'in' expression body
iterator       := 'ref'? ('mut' | 'dmut')? (Identifier | '_')
                  (':' type)?
body           := block
                | expression ';'?
\end{lstlisting}

As for a conditional, the body can be a single expression.

\subsection{Evaluation}

The iterated value is evaluated once, before the first iteration. Changing a
variable that it was computed from does not change the number of iterations. The
body is then evaluated for each element, in order, with the iterators bound to
it. \token{break} and \token{continue} behave as in a \token{loop}
(Section~\ref{sec:flow:inf_loop}); after a \token{continue}, the iterators move
to the next element.

Throughout this section, the iterator that receives an element is called the
\textit{value iterator}, and the one that receives its position the
\textit{index iterator}.

\begin{lstlisting}[style=coloredverbatim, escapechar=@]
let a = copy [1, 2, 3];

for i, v in a {
  println("Value at index : ", i , " equals to : ", v);
}
\end{lstlisting}

\subsection{Iterators}

The number of iterators, and what each receives, depends on the type of the
iterated value. Any other number of iterators is an error, E4244.

\begin{center}
  \begin{adjustbox}{max width=\linewidth}
    \begin{tabular}{|l|l|l|}
      \hline
      Iterated value & One iterator & Two iterators\\[0pt]
      \hline
      \hline
      range \token{..T} & the value, of type \token{T} & --\\[0pt]
      array \token{[T; N]}, slice \token{[T]} & the element, of type \token{T} & the index, of type \token{usize}, then the element\\[0pt]
      tuple & the field & the index, of type \token{usize}, then the field\\[0pt]
      map \token{[K => V]} & the key & the key, then the value\\[0pt]
      class & \multicolumn{2}{l|}{as defined by the class (Section~\ref{sec:class_override_for_loop})}\\[0pt]
      \hline
    \end{tabular}
  \end{adjustbox}
\end{center}

\begin{lstlisting}[style=coloredverbatim, escapechar=@]
for i in 0 .. 10 {
  println("Index iteration : ", i);
}

for v in [1, 2, 3] {
  println("Value iteration : ", v);
}

for i, v in [1, 2, 3] {
  println("Value and index iteration : ", i, ' ', v);
}
\end{lstlisting}

A string is a slice of characters, so a \token{for} over a \token{[c8]} visits
its code units, not its characters (Section~\ref{sec:types:chars_strings}).

The type of an iterator is inferred. It can be written after \token{:}, but must
then be the type that would have been inferred; the explicit type exists to
carry mutability (Section~\ref{sec:flow:mut_value_iterators}).

\begin{lstlisting}[style=coloredverbatim, escapechar=@]
let a = copy [1, 2, 3];

for i: usize, v: i32 in a {
  println("Value at index : ", i , " equals to : ", v);
}
\end{lstlisting}

An iterator is not an lvalue: the body cannot assign it (E4153), except through
a mutable reference iterator (Section~\ref{sec:flow:mut_ref_iterators}). The
index iterator is always immutable, and takes no decorator.

An iterator is a variable, subject to the rule that a variable must be used
(Section~\ref{sec:memory:unused_variables}). The wildcard \token{\_} in place of
its name declares nothing, and is the way to ignore an element or an index:
\token{for \_ in 0 .. n} repeats \token{n} times, and \token{for \_, v in a}
visits the values of \token{a}.

\subsection{Iteration over a tuple}

The fields of a tuple can have different types, so no single iterator type fits
them all. A \token{for} over a tuple is therefore unrolled at compile time: the
body is copied once per field, with the value iterator of the type of that field.
The loop of line~2 in Listing~\ref{lst:flow:ref_value_iter_ex} becomes the
M-YIL of Listing~\ref{lst:flow:copy_value_iter_tu}, in which each copy of the
body is a block of its own. \token{break} and \token{continue} are allowed in
the body, and jump out of the unrolled sequence, or to the next copy, as they
would in any loop.

\subsection{Mutable value iterator}
\label{sec:flow:mut_value_iterators}

A value iterator is immutable by default. It can be declared \token{mut} or
\token{dmut}, but only when it borrows mutable data from the iterated value,
which must then be given with \token{alias}. Because iterators are not lvalues,
this mutability lets the body modify the data the iterator borrows, never the
iterator itself: an iterator behaves like a parameter
(Section~\ref{sec:mutable_parameter}), which cannot be reassigned unless it is a
reference.

\begin{lstlisting}[style=coloredverbatimError, escapechar=@]
let dmut a = dcopy [[1, 2], [3, 4]];
for dmut v in alias a {
  v[0] = 10; // ok, aliasing /a/
  @\hb{v = copy [8, 9];}@ // error, /v/ is not a lvalue
}

println(a);
\end{lstlisting}

An iterator that borrows no mutable data cannot be mutable (E4137): its
mutability could have no effect on the iterated value.

\begin{lstlisting}[style=coloredverbatimError, escapechar=@, caption=Useless mutable iterator, label=lst:flow:useless_mut_iterator]
let dmut a = copy [1, 2, 3];
for @\hb{mut}@ v in alias a { // error, /v/ does not borrow mutable data
  @\hb{v = 1;}@ // would have no effect on /a/
}
\end{lstlisting}

\subsection{Reference iterators}

\token{ref} on the value iterator binds it to each element by reference, instead
of copying the element into it at each iteration. The iterated value must be an
lvalue. It need not be mutable, as long as the iterator is not
(Section~\ref{sec:flow:mut_ref_iterators}). This saves the copy of large
elements; in the next example, both loops print the same thing, but the second
copies a pointer at each iteration where the first copies 1024 integers.

\begin{lstlisting}[style=coloredverbatim, escapechar=@]
let a: [[i32 ; 1024]] = copy [[0 ; 1024], [1 ; 1024]];
for v in a { // copy 1024 i32, at each iteration
  println(v);
}

for ref v in a { // reference, copying only 1 pointer (8 bytes)
  println(v);
}
\end{lstlisting}

Reference iterators apply to tuples, arrays and slices. In
Listing~\ref{lst:flow:ref_value_iter_ex}, the loop of line~2 copies the fields of
the tuple, and the loop of line~6 references them; the loops of lines~11 and~15
do the same on an array. The M-YIL of the first two loops is
Listing~\ref{lst:flow:copy_value_iter_tu} and
Listing~\ref{lst:flow:ref_value_iter_tu}.

\begin{lstlisting}[style=coloredverbatim, caption=iteration using reference value iterator, label=lst:flow:ref_value_iter_ex]
let a = ([0 ; 100], [1 ; 10]);
for v in a {
  println(v);
}

for ref v in a {
  println(v);
}

let b: [[i32 ; 1024] ; 2] = [[0 ; 1024], [1 ; 1024]];
for v in b {
  println(v);
}

for ref v in b {
  println(v);
}
\end{lstlisting}

\begin{lstlisting}[style=myilVerb, caption=By copy iteration on a tuple, label=lst:flow:copy_value_iter_tu]
cte for {
  {
    let v(#48) : [i32 ; 100us] = [0 ; 100us];
    {
      std::io::println!{[i32 ; 100us]}::println (v(#48));
      <unit-value>
    }
  };
  {
    let v(#s7) : [i32 ; 10us] = [1 ; 10us];
    {
      std::io::println!{[i32 ; 10us]}::println (v(#s7));
      <unit-value>
    }
  }
};
\end{lstlisting}

\begin{lstlisting}[style=myilVerb, caption=By reference iteration on a tuple, label=lst:flow:ref_value_iter_tu]
cte for {
  {
    let ref v(#1gd) : [i32 ; 100us] = ref (a(#3t).0);
    {
      std::io::println!{[i32 ; 100us]}::println (v(#1gd));
      <unit-value>
    }
  };
  {
    let ref v(#1hd) : [i32 ; 10us] = ref (a(#3t).1);
    {
      std::io::println!{[i32 ; 10us]}::println (v(#1hd));
      <unit-value>
    }
  }
};
\end{lstlisting}

\subsection{Mutable reference iterators}
\label{sec:flow:mut_ref_iterators}

A reference iterator declared \token{mut} or \token{dmut} is an lvalue:
assigning it assigns the element it references. It follows the rules of any
mutable reference variable (Section~\ref{sec:mut_ret_param}). The iterated value
must have the mutability of the iterator, and the contract is accepted by writing
\token{ref} in front of the iterated value.

\begin{lstlisting}[style=coloredverbatim, escapechar=@]
let dmut a = copy [1, 2, 3];

for i, ref mut v in ref a {
  //                ^^^ mandatory
  v = cast!i32(i) + 9;
}

println(a); // [9, 10, 11]
\end{lstlisting}

The mutability of the iterator must not exceed that of the iterated value
(E4045). In the next example, \token{a} can modify its own elements, the inner
slices, but not the integers inside them. The loop of line~4 asks for a deeply
mutable iterator, which could modify those integers, and is rejected. The loop
of line~9 asks for one level of mutability, which \token{a} has, and replaces
the inner slices without touching \token{y}.

\begin{lstlisting}[style=coloredverbatimError, escapechar=@]
let y = copy [3, 4];
let mut a: [mut [i32]] = dcopy [[1, 2], y];

for @\hb{ref dmut}@ v in ref a {
  //    ^^^, error, incompatible mut [i32] -> mut [mut i32]
  @\hb{v[0] = 9;}@ // would modify values that /a/ does not own
}

for ref mut v in ref a {
  //    ^^^ Ok, slice in /a/ is mutable
  v = copy [10, 11];
}

println(a); // [[10, 11], [10, 11]]
println(y); // [3, 4], untouched as guaranteed
\end{lstlisting}

\subsection{Value and type}

A \token{for} has no value, and its type is \token{void}, for the reason given
for \token{while} (Section~\ref{sec:flow:while_value}): the loop can end by
running out of elements. A \token{break} in a \token{for} carries no value.

\subsection{Compile-time iteration}
\label{sec:flow:cte_for}

\token{cte for} unrolls the loop at compile time, as a loop over a tuple is. The
iterated value must be known at compile time (E4256), and the iterators are then
compile-time values in each copy of the body. A test on an iterator is
therefore known at compile time, and is written \token{cte if}
(Section~\ref{sec:flow:cte_if}): a plain \token{if} on it is rejected, as any
test whose outcome the compiler knows (Section~\ref{sec:flow:cte}).

\begin{lstlisting}[style=coloredverbatim, caption=A compile-time loop]
fn main() {
  cte for i in 0 .. 4 {
    cte if i == 2 {
      println("two");
    } else {
      println(i);
    }
  }
}
\end{lstlisting}

\subsection{Diagnostics}

\begin{center}
  \begin{adjustbox}{max width=\linewidth}
    \begin{tabular}{|l|l|}
      \hline
      Code & Raised when\\[0pt]
      \hline
      \hline
      E4244 & the number of iterators does not fit the iterated value\\[0pt]
      E4153 & the body assigns an iterator that is not a mutable reference\\[0pt]
      E4137 & a mutable iterator borrows no mutable data\\[0pt]
      E4045 & the mutability of an iterator exceeds that of the iterated value\\[0pt]
      E4095 & a \token{break} carries a value\\[0pt]
      E4256 & the value iterated by a \token{cte for} is not known at compile time\\[0pt]
      \hline
    \end{tabular}
  \end{adjustbox}
\end{center}
